% !TeX root = ../papers/nc-p3.tex
\section{導手直線に関する残余スキームと非被約境界の排除}
\label{sec:boundary-conductor-residual}
\begin{theorem}\label{thm:boundary-conductor-residual}
四辺有限な標準有限Frobenius窓で、整な八余因子八次曲面の節点・尖点核の境界Hankel元は存在しない。
導手直線に関して一度残余を取ると、導手に支えを持つ部分も含めて前節の有限退化上に写る。
\end{theorem}
\begin{proof}
入力点スキームを $X$ とする。その各局所Artin環 $R$ で入力の一座標を1にし、
$N=y(z)y(w)$ と書く。階数1で少なくとも一つの対角成分が単元であり、その平方根が取れる。
$e=y(p)$ と置く。$p$ 行の係数イデアルは $(e)$ である。一方この行は節点で
$a(z)(a_3+a_0z)$、尖点で $a(z)(a_0-a_1z)$ に等しい。$a$ は剰余体上でも非零で
$a$ による乗法 $S_1\to S_4$ は分裂単射である。従って $(e)$ は導手直線 $C$ の
$X$ 上のイデアルと一致する。$R'=R/\operatorname{ann}(e)$ と置く。$R'=0$ の箇所は除ける。
$R'$ 上で $s=a_0/e$ とし、節点で $t=a_3/e$、尖点で $t=a_1/e$ とする。
割算はイデアル $(e)\simeq R'$ によるもので選択に依らない。この二元は単位イデアルを生成する。
$p$ 行を一度 $e$ で割ると $R'$ 上で
\[
 y(z)=a(z)(t+sz)\quad\text{または}\quad y(z)=a(z)(s-tz)
\]
となる。$z=0$ を入れれば、節点で $e(st-1)=0$、尖点で $e(s^2-1)=0$ である。
以下で $R'$ 自体の上の等式 $st=1$ または $s^2=1$ を示す。

まず $C$ 上の台点で剰余体に落とす。$a=g h$ と書け、$h$ は非零二項一次式である。
節点では $g=z$、尖点では $g=z^2$ を二項二次式として読む。$N=yy^{\mathsf t}$ から特解を引き
$g(z)g(w)$ で割ると $B(a,Ta)$ は順に
\[
 h(z)h(w)\{t^2+(st-1)(z+w)+s^2zw\},\qquad
 h(z)h(w)\{s^2-1-st(z+w)+t^2zw\}
\]
となる。一般に $B(gh,b)$ が $h(z)h(w)$ で割れれば $h\mid b$ である。
$h$ の根で片方の変数を評価するだけでよく、無限遠も同次形で同様である。
$h$ を除くと節点の $B(z,S_2)$ は $1,zw$ の包、尖点の $B(z^2,S_2)$ は $z+w,zw$ の包である。
従って剰余体でそれぞれ $st=1,s^2=1$、特に $s$ は単元であり、節点では $t$ も単元となる。
$C$ の外では $e$ は初めから単元なので同じ結論が成立する。

$u=t/s$ とし、$q=st-1$ または $q=s^2-1$ とする。節点で $u$ は単元、尖点で $u$ は有限である。
前節のSchur恒等式は $a_3=u a_0$ または $a_1=u a_0$ の超平面束上の恒等式で、
極以外では $a_0=0$ でも有効である。$y(z)y(w)$ との差を取ると次の等式を得る。
\[
 B(a,Ta-U_u(a))=
 \begin{cases}
 (q/u) h(z)h(w)(u+z)(u+w)&\text{節点},\\
 q h(z)h(w)(1-uz)(1-uw)&\text{尖点}.
 \end{cases}
\]
ここで $h=a_1+a_2z$ または $a_2+a_3z$ と取る。実際 $eq=0$ なので
$q a(z)=q g(z)h(z)$ であり、両変数でこの置換ができる。
$B(a,S_3)$ は随所で階数3の直和因子である。$C$ 上の剰余体で右辺の $q$ 以外の因子は
この直和因子に入らない。節点では $h(z)h(w)$ を除いた後の $z+w$ の係数が1、尖点では定数項が1で、
先に求めた二空間から外れるためである。商のある座標が単元なので $q=0$ が $R'$ で従う。

これで $B(a,Ta-U_u(a))=0$、従って $Ta-U_u(a)=\lambda a$ が $R'$ 上で成立する。
$s,t$ が同時に符号を変えても $u$ は同じで、$\lambda$ も原始列 $a$ から一意に決まる。
従って $X'$ を $\mathcal O_X/\operatorname{ann}(\mathcal I_C\mathcal O_X)$ で定めると、
$X'$ 全体が前節の完備化した固有直線入射に入る。各スカラーファイバーは入力点または直線の
$X'$ との切断の閉部分であり、相対二次式の全射と前節の余核長16の議論がそのまま使える。
残る商は $X\cap C$ で長さ3以下である。従って
\[
            \ell(X)=\ell(X')+\ell(X\cap C)\le16+3=19 ,
\]
所要長20に反する。特に $e^2=0$ の局所成分や導手の両端の重根もこの議論に含まれる。
\end{proof}

\path{computations/verify_boundary_conductor_residual.py} は節点・尖点のSchur恒等式を
超平面束上で記号検算し、導手での二つのBezout像を係数比較する。
